% GENERATED FILE. Edit canonical lessons, then run npm run generate:books.
% canonical-source-hash: fnv1a64:4522df92257600c0
% canonical-lessons: FR-C131-occupe, FR-C131-libre, FR-C131-ouvert, FR-C131-ferme, FR-C131-calme

\chapter{Busy, Free, Open, Closed, Calm}
\label{ch:fr-busy-free-open-closed-calm}
\hlchaptermodality{131}

\begin{chapteropening}
\textbf{By the end of this chapter:} \emph{I can name busy, free, open, closed, calm.}

Complete the last lesson of chapter 131: i can name busy, free, open, closed, calm.
\end{chapteropening}

\section[occupé]{occupé — busy}

\label{lesson:FR-C131-occupe}



\begin{quote}
Before the new one: say the French for brown, then the French for pretty.
\end{quote}

\subsection*{You'll want to know: occupé}
\textbf{occupé} — \textquotedblleft{}busy\textquotedblright{}.

\begin{grammarlens}[title={What you've built}]
One more word for this chapter.
\end{grammarlens}

\subsection*{Your turn}
\begin{itemize}
\raggedright
  \item \emph{Say it:} \emph{occupé}
  \item \emph{Say it:} \emph{occupé}, once more
  \item \emph{Say it:} say \emph{occupé}
  \item \emph{Recall:} say the French for brown, then the French for pretty, then say \emph{occupé} again
\end{itemize}

\begin{tcolorbox}[breakable,colback=teal!4,colframe=teal!35!black,title={Before you move on}]
What does occupé mean? (\textquotedblleft{}Busy\textquotedblright{}.) Say it once more.
\end{tcolorbox}
\section[libre]{libre — free}

\label{lesson:FR-C131-libre}



\begin{quote}
Before the new one: say the French for pretty, then the French for busy.
\end{quote}

\subsection*{You'll want to know: libre}
\textbf{libre} — \textquotedblleft{}free\textquotedblright{}.

\begin{grammarlens}[title={What you've built}]
One more word for this chapter.
\end{grammarlens}

\subsection*{Your turn}
\begin{itemize}
\raggedright
  \item \emph{Say it:} \emph{libre}
  \item \emph{Say it:} \emph{libre}, once more
  \item \emph{Say it:} say \emph{libre}
  \item \emph{Recall:} say the French for pretty, then the French for busy, then say \emph{libre} again
\end{itemize}

\begin{tcolorbox}[breakable,colback=teal!4,colframe=teal!35!black,title={Before you move on}]
What does libre mean? (\textquotedblleft{}Free\textquotedblright{}.) Say it once more.
\end{tcolorbox}
\section[ouvert]{ouvert — open}

\label{lesson:FR-C131-ouvert}



\begin{quote}
Before the new one: say the French for busy, then the French for free.

\textbf{Ouvert} belongs to \textbf{ouvrir}, to open, which takes the endings of a verb in -er: \emph{j'ouvre}. Say \emph{j'ouvre}, then its opposite, \emph{je ferme}.
\end{quote}

\subsection*{You'll want to know: ouvert}
\textbf{ouvert} — \textquotedblleft{}open\textquotedblright{}.

\begin{grammarlens}[title={What you've built}]
One more word for this chapter.
\end{grammarlens}

\subsection*{Your turn}
\begin{itemize}
\raggedright
  \item \emph{Say it:} \emph{ouvert}
  \item \emph{Say it:} \emph{ouvert}, once more
  \item \emph{Say it:} say \emph{ouvert}
  \item \emph{Recall:} say the French for busy, then the French for free, then say \emph{ouvert} again
\end{itemize}

\begin{tcolorbox}[breakable,colback=teal!4,colframe=teal!35!black,title={Before you move on}]
What does ouvert mean? (\textquotedblleft{}Open\textquotedblright{}.) Say it once more.
\end{tcolorbox}
\section[fermé]{fermé — closed}

\label{lesson:FR-C131-ferme}



\begin{quote}
Before the new one: say the French for free, then the French for open.

\textbf{Fermé} is the closed form of \textbf{fermer}, to close. Say \emph{fermer}, then \emph{fermé}.
\end{quote}

\subsection*{You'll want to know: fermé}
\textbf{fermé} — \textquotedblleft{}closed\textquotedblright{}.

\begin{grammarlens}[title={What you've built}]
One more word for this chapter.
\end{grammarlens}

\subsection*{Your turn}
\begin{itemize}
\raggedright
  \item \emph{Say it:} \emph{fermé}
  \item \emph{Say it:} \emph{fermé}, once more
  \item \emph{Say it:} say \emph{fermé}
  \item \emph{Recall:} say the French for free, then the French for open, then say \emph{fermé} again
\end{itemize}

\begin{tcolorbox}[breakable,colback=teal!4,colframe=teal!35!black,title={Before you move on}]
What does fermé mean? (\textquotedblleft{}Closed\textquotedblright{}.) Say it once more.
\end{tcolorbox}
\section[calme]{calme — calm, quiet}

\label{lesson:FR-C131-calme}



\begin{quote}
Before the new one: say the French for open, then the French for closed.
\end{quote}

\subsection*{You'll want to know: calme}
\textbf{calme} — \textquotedblleft{}calm, quiet\textquotedblright{}.

\begin{grammarlens}[title={What you've built}]
One more word for this chapter.
\end{grammarlens}

\subsection*{Your turn}
\begin{itemize}
\raggedright
  \item \emph{Say it:} \emph{calme}
  \item \emph{Say it:} \emph{calme}, once more
  \item \emph{Say it:} say \emph{calme}
  \item \emph{Recall:} say the French for open, then the French for closed, then say \emph{calme} again
\end{itemize}

\begin{tcolorbox}[breakable,colback=teal!4,colframe=teal!35!black,title={Before you move on}]
What does calme mean? (\textquotedblleft{}Calm, quiet\textquotedblright{}.) Say it once more.
\end{tcolorbox}
